\subsection{PROPERTIES PECULIAR TO OPEN MAPPINGS}

PROPOSITION 5. Let X, Y be two topological spaces, f a mapping of X into Y, B a base of the topology of X. Then the following statements are equivalent:

\begin{enumerate}
\def\labelenumi{\alph{enumi})}
\item
  \(f\) is an open mapping.
\item
  For each \(\mathbf{U} \in \mathfrak{B}\) , \(f(\mathbf{U})\) is open in \(\mathbf{Y}\) .
\item
  For each \(x \in \mathbf{X}\) and each neighbourhood \(\mathbf{V}\) of \(x\) in \(\mathbf{X}\) , \(f(\mathbf{V})\) is a neighbourhood of \(f(x)\) in \(\mathbf{Y}\) .
\end{enumerate}

The equivalence of a) and b) follows immediately from the definitions and from \((\mathrm{O}_{\mathrm{I}})\) ; the equivalence of a) and c) is a consequence of Proposition 1 of §1, no. 2.

PROPOSITION 6. Let R be an equivalence relation on a topological space X; then the following three conditions are equivalent:

\begin{enumerate}
\def\labelenumi{\alph{enumi})}
\item
  The relation \(\mathbf{R}\) is open.
\item
  The interior of each subset which is saturated with respect to R is saturated with respect to R.
\item
  The closure of each subset which is saturated with respect to R is saturated with respect to R.
\end{enumerate}

By taking complements (§ 1, no. 6, formula (2)) we see that b) and c) are equivalent. Let us show that b) implies a): suppose condition b) is satisfied and let U be an open subset of X, V its saturation with respect to R; then \(\dot{V} \supset U\) , and since by hypothesis \(\dot{V}\) is saturated, it follows that \(\dot{V} = V\) , and therefore the saturation of U is open. Conversely, suppose condition a) is satisfied, and let A be a saturated set; if B is the saturation of \(\mathring{A}\) , then \(\mathring{A} \subset B \subset A\) , and since B is open by hypothesis it follows that \(B = \mathring{A}\) .

PROPOSITION 7. Let R be an open equivalence relation on a topological space X, and let \(\varphi: \mathbf{X} \to \mathbf{X}/\mathbf{R}\) be the canonical mapping. If A is any subset of X which is saturated with respect to R, then the closure (resp. the interior) of \(\varphi(\mathbf{A})\) in X/R is \(\varphi(\overline{\mathbf{A}})\) (resp. \(\varphi(\mathring{\mathbf{A}})\) ).

Each of the two assertions of the proposition can be deduced from the other by taking complements and using formula (2) of § 1, no. 6 and the fact that if B is a saturated subset of X then \(\varphi(\complement B) = \complement \varphi(B)\) . By virtue of Proposition 6, \(\overline{A}\) is saturated; hence \(\varphi(\overline{A})\) is closed in \(X/R\) , and since \(A \subset \overline{A}\) we have \(\overline{\varphi(A)} \subset \varphi(\overline{A})\) , so that \(\varphi(A) \subset \varphi(\overline{A})\) . But since \(\varphi\) is continuous, \(\varphi(\overline{A}) \subset \overline{\varphi(A)}\) ( \(\S 2\) , no. 1, Theorem 1), and the result follows.

PROPOSITION 8. Let \((\mathbf{X}_t)_{t \in \mathbf{I}}\) , \((\mathbf{Y}_t)_{t \in \mathbf{I}}\) be two families of topological spaces indexed by the same set I. For each \(t \in \mathbf{I}\) let \(f_t\) be an open mapping of \(\mathbf{X}_t\) into \(\mathbf{Y}_t\) , and suppose that \(f_t\) is surjective for all but a finite number of indices. Then the product mapping \(f: (x_t) \to (f_t(x_t))\) of \(\prod_{t \in \mathbf{I}} \mathbf{X}_t\) into \(\prod_{t \in \mathbf{I}} \mathbf{Y}_t\) is open.

By virtue of Proposition 5 we need only prove that the image under \(f\) of any elementary set \(\prod_{t \in I} A_t\) in \(\prod_{t \in I} X_t\) is open in \(\prod_{t \in I} Y_t\) . But this image is \(\prod_{t \in I} f_t(A_t)\) , and the hypotheses imply that \(f_t(A_t)\) is open in \(Y_t\) for each \(t \in I\) , and that \(f_t(A_t) = Y_t\) for all but a finite number of indices; whence the result.

COROLLARY. Let \((\mathbf{X}_t)_{t \in I}\) be a family of topological spaces, and for each \(t \in I\) let \(\mathbf{R}_t\) be an equivalence relation on \(\mathbf{X}_t\) , and let \(f_t\) be the canonical mapping \(\mathbf{X}_t \to \mathbf{X}_t / \mathbf{R}_t\) . Let \(\mathbf{R}\) be the equivalence relation in \(\mathbf{X} = \prod_{t \in I} \mathbf{X}_t\) ``for each \(t \in I\) , \(\operatorname{pr}_t(x) = \operatorname{pr}_t(y) (\bmod \mathbf{R}_t)"\)

between \(x\) and \(y\) , and let \(f\) be the product mapping \((x_{t}) \to (f_{t}(x_{t}))\) of \(X\) into \(\prod_{t \in I} (X_{t}/R_{t})\) . If each of the relations \(R_{t}\) is open, then the relation \(R\) is open, and the bijection associated with \(f\) is a homeomorphism of \(X/R\) onto \(\prod_{t \in I} (X_{t}/R_{t})\) .

R is the relation \(f(x)=f(y)\) . Since f is continuous and open by Proposition 8 above and § 4, no. 1, Corollary 1 of Proposition 1, the result follows from Proposition 3 of no. 2.

In particular, if R (resp. S) is an open equivalence relation on a topological space X (resp. Y), then the canonical bijection of

\[
(\mathbf {X} \times \mathbf {Y}) / (\mathbf {R} \times \mathbf {S}) \quad \text { onto } \quad (\mathbf {X} / \mathbf {R}) \times (\mathbf {Y} / \mathbf {S})
\]

\pandocbounded{\includegraphics[keepaspectratio]{images/part001_bebaa45c.jpg}}

is a homeomorphism. If R and S are not assumed to be open, this bijection is continuous but is not necessarily a homeomorphism, even if one of the relations R, S is the relation of equality (Exercise 6).

